\chapter*{Streszczenie}
%\addcontentsline{toc}{chapter}{Streszczenie}  
Niniejsza praca prezentuje opracowanie aplikacji mobilnej typu PWA, która umożliwia interaktywną rozgrywkę polegającą na odgadywaniu lokalizacji na podstawie zdjęć oraz weryfikowaniu obecności w wybranych miejscach z wykorzystaniem geolokalizacji. Praca rozpoczyna się od analizy sytuacji na rynku i uzasadnienia wyboru tematu, a następnie przedstawia szczegółowy podział prac oraz wprowadzenie do dziedziny aplikacji mobilnych i gier edukacyjno-rozrywkowych.

W pracy omówiono technologie i narzędzia wykorzystane w projekcie, w tym frameworki frontendowe i backendowe, systemy mapowe oraz mechanizmy geolokalizacji. Na podstawie analizy wymagań użytkowych i systemowych zaprojektowano architekturę aplikacji, uwzględniając zarówno aspekt funkcjonalny, jak i niefunkcjonalny, w tym responsywność, estetykę interfejsu i wydajność.

Część implementacyjna obejmuje szczegółowy opis tworzenia frontendu i backendu, integrację z usługami zewnętrznymi oraz implementację mechanizmów rozgrywki, takich jak tryby klasyczny i geolokalizacyjny, system punktacji, historia rozgrywek i ranking użytkowników. W pracy przedstawiono także konfigurację aplikacji PWA, zarządzanie sesjami użytkowników oraz obsługę zdjęć lokalizacji.

Rozdział dotyczący testowania opisuje przeprowadzone testy manualne funkcjonalności oraz weryfikację stabilności i poprawności działania aplikacji na różnych urządzeniach i przeglądarkach. W podsumowaniu pracy omówiono osiągnięte cele oraz wskazano potencjalne kierunki dalszego rozwoju projektu, takie jak rozbudowa bazy lokalizacji, wprowadzenie nowych trybów gry czy integracja z dodatkowymi usługami sieciowymi.

Praca przedstawia pełen proces projektowania, implementacji i testowania aplikacji mobilnej typu PWA, łączącej funkcjonalności edukacyjne i rozrywkowe w intuicyjnym i atrakcyjnym interfejsie użytkownika.

\chapter*{Abstract}
%\addcontentsline{toc}{chapter}{Abstract}
This thesis presents the development of a mobile PWA application that enables interactive gameplay based on guessing locations from photos and verifying presence in selected places using geolocation. The work begins with a market analysis and justification for the topic choice, followed by a detailed work distribution and an introduction to the field of mobile applications and educational-entertainment games.

The thesis discusses the technologies and tools used in the project, including frontend and backend frameworks, mapping systems, and geolocation mechanisms. Based on the analysis of user and system requirements, the application architecture was designed, considering both functional and non-functional aspects, including responsiveness, interface aesthetics, and performance.

The implementation section includes a detailed description of frontend and backend development, integration with external services, and the implementation of gameplay mechanisms such as classic and geolocation modes, scoring system, game history, and user ranking. The thesis also presents the configuration of the PWA application, user session management, and photo handling for locations.

The testing chapter describes the manual functional tests conducted and the verification of the application's stability and correctness across different devices and browsers. In the conclusion, the achieved goals are discussed, and potential directions for further project development are indicated, such as expanding the location database, introducing new game modes, or integrating with additional web services.

The thesis presents the complete process of designing, implementing, and testing a mobile PWA application that combines educational and entertainment functionalities in an intuitive and attractive user interface.
